% !TeX root = ../main.tex
\chapter{Week 3 Assignment, Part B in re: Child Labor}
\label{ch:child-labor}
\section*{Scenario}
Consider the following scenario.
Your design is a high-profile product for your company in the U.S.A. As such, the media
makes every effort to monitor the project's progress. Several vendors have been
selected for sourcing the parts, all overseas. Things are looking good until your
manager comes into your office and tells you that the media has just released a story
that the manufacturer of your highest-cost part uses child labor, that is children as
young as 14. The legal working age in the manufacturer's country is 14. To complicate
matters, to change vendors would result in an additional cost of 15% for this high-cost
item. You have been invited to a meeting to discuss how to proceed. Using your text
from your ethics course (and the internet as well if you wish), prepare a
recommendation for the meeting. This recommendation should not be based on opinion
but rather on problem-solving approaches when dealing with cross-cultural ethics
problems. Would your recommendation be different, and if so, why, if the manufacturer
used children as young as 12?

\section*{Recommendation}
In such a scenario, I would direct the company to investigate the supplier before they jump to any rash decisions. Cross-cultural ethics issues shouldn’t be viewed only from the lens of US customs, or only by the legality in the foreign country. The company needs to take all conditions into account – laws, working conditions, safety and education, and the effects that the decision can have on all stakeholders.
Workers who are 14 may be ethically defensible if the work the 14 year olds are doing is legal, safe, voluntary, and pay is commensurate with the work, while not interfering with their education. Further, the International Labour Organization does allow workers as young as 14 for some countries, so long as the work is deemed non-hazardous (International Labour Organization, 1973). The company must verify the allegations and apply reasonable labor standards, rather than changing suppliers at the whim of social or media pressures.
The additional 15\% cost in changing suppliers is worth considering – but it alone shouldn’t be the deciding factor if workers are indeed being exploited or put at risk. If the investigation finds issues such as hazardous working conditions, excessive hours that interfere with education requirements, coercion of workers, or the like, the company should require corrective action and move to another supplier if problems are not resolved.
Workers that are as young as 12 would shift my recommendation. Since the scenario states that the working age is 14 in this hypothetical, a child at 12 would be violating the labor laws in the country. If this were discovered, I would demand suspension of new purchases pending investigation and remediation of the issue. Ethically, this would force the use of a new supplier despite the cost difference, if it were found to be true and it weren’t resolved.
